\section*{Chapter \ref{ch:rc:modelling-overnight-rate-products} --- Modelling Overnight-Rate Products}

\begin{solution}{exo:rc:modelling-overnight-rate-products:1}
An interbank rate for the current period was fixed at its start, so its caplet had
a known payoff. A compounded overnight rate for the current quarter is known only at
the quarter's end: the fixings still to come make the caplet a genuine option.
\end{solution}

\begin{solution}{exo:rc:modelling-overnight-rate-products:2}
By $\sigma^2\delta/3 = \sigma^2\times0.0833$: total variance $\sigma^2\times2.0833$ against
$\sigma^2\times2$, 4.2\% more.
\end{solution}

\begin{solution}{exo:rc:modelling-overnight-rate-products:3}
27 October 2027 is 1.077 years from spot and 8 December 1.192; the quarter ends at
1.25: $(1.25-1.077)/0.25 = 69\%$ and $(1.25-1.192)/0.25 = 23\%$.
\end{solution}

\begin{solution}{exo:rc:modelling-overnight-rate-products:4}
$\frac{\delta^3/3}{\delta^2S+\delta^3/3} = \frac{\delta/3}{S+\delta/3} = \frac{0.0833}{0.3333} = 25\%$ for
$S=0.25$, as in the text (25.2\% with mean reversion).
\end{solution}

\begin{solution}{exo:rc:modelling-overnight-rate-products:5}
On the meeting days: its remaining uncertainty falls in steps on 27 October and 8
December, and after 8 December it has none. Its theta is concentrated on those days
(small in between), unlike the smooth decay of the generalised model; a trader must
hedge and mark it meeting by meeting.
\end{solution}

\begin{solution}{exo:rc:modelling-overnight-rate-products:6}
A term rate is computed from derivatives on the overnight rate; if derivatives on the
term rate became large, it would rest partly on trading in itself and could be moved
by the flows it is meant to measure. Restricting its derivatives to end users'
hedges keeps the liquidity in overnight products. A desk that sells term-rate caps
hedges with overnight-rate options and is left with the in-period variance and the
difference between the published term rate and the realised compounded rate.
\end{solution}

\begin{solution}{exo:rc:modelling-overnight-rate-products:7}
The first caplet, USD~13\,623, exists only in the compounded cap; the other seven
caplets are worth USD~24\,599 more in total than their forward-looking twins, from
their in-period variance: 13\,623 plus 24\,599 is the 38\,222.
\end{solution}

\begin{solution}{exo:rc:modelling-overnight-rate-products:8}
Both assumptions are wrong for compounded rates: the current period's caplet is
still an option and must be priced, and each caplet's variance runs to the end of
its period, not its start. The pricer undervalues every cap, here by USD~38\,222 on
USD~100 million over two years, and its vega is on the wrong dates.
\end{solution}

\begin{solution}{pb:rc:modelling-overnight-rate-products:1}
\textbf{1.} USD~288\,963 on compounded SOFR, USD~250\,741 on the term rate.
\textbf{2.} It includes the current quarter's caplet and each caplet carries risk
until the end of its period.
\textbf{3.} USD~13\,623.
\textbf{4.} The compounded-SOFR cap, since the loan pays compounded SOFR; the term-rate
cap leaves the difference between the loan's rate and the term rate.
\textbf{5.} 25.2\% and 4.8\%.
\textbf{6.} USD~40\,944 by the closed form; the simulation agrees to within a fifth of a
standard error.
\textbf{7.} The time grid missed the period's dates (312.5 steps truncated to 312), so the
simulated period was shorter than the quarter and the rate too low: a 6\% bias with a
tight standard error.
\textbf{8.} A variance larger by $\delta/(3S) = 8.3\%$; Hull--White gives an in-period
share of 7.9\%, i.e. 8.6\% more, almost the same.
\textbf{9.} Hull--White's in-period variance is close to $\sigma^2\delta^3/3$, exactly what a
linear decay of the rate's volatility produces.
\textbf{10.} The volatility level and, for this out-of-the-money strike, the forward
rates.
\textbf{11.} 27 October and 8 December 2027.
\textbf{12.} After the 8 December decision, three weeks before the quarter ends.
\textbf{13.} Zero: nothing uncertain remains.
\textbf{14.} With futures on the meeting months and options on them, sized to each
meeting's weight in the quarter.
\textbf{15.} The risk: when theta is earned and where the vega sits.
\textbf{16.} The compounded-SOFR cap, which matches the loan; the extra USD~38\,222 buys
protection on the current quarter and removes basis risk.
\textbf{17.} The difference between term SOFR and compounded SOFR over each quarter,
including the in-period variance it cannot hedge in term-rate options.
\textbf{18.} Both limit the loan's interest; the term-rate cap leaves a small basis they
tolerate.
\textbf{19.} Named result: \emph{the SOFR cap}: USD~288\,963 on compounded SOFR against
USD~250\,741 on the term rate; the USD~38\,222 gap is the current quarter's caplet
(13\,623) and in-period variance (24\,599).
\textbf{20.} Risk that keeps accruing until the end of the period.
\end{solution}

\begin{solution}{iq:rc:modelling-overnight-rate-products:1}
Euribor fixes at the start of the period, so its caplet's risk ends then;
compounded €STR fixes at the end, so its caplet carries risk through the accrual
period, is worth more, and the current period's caplet is an option too.
\iqlookfor{the timing of the fixing and its consequence for variance.}
\end{solution}

\begin{solution}{iq:rc:modelling-overnight-rate-products:2}
Model the forward of the \omterm{def:rc:modelling-overnight-rate-products:rates}{backward-looking rate} for each period up to the period's
end, with a volatility that decays to zero during the period (the \omterm{def:rc:modelling-overnight-rate-products:gfmm}{generalised
forward market model}); before the period it equals the forward-looking forward, so
one process gives both kinds of rate.
\iqlookfor{forwards alive to the end of the period, and decaying volatility.}
\end{solution}

\begin{solution}{iq:rc:modelling-overnight-rate-products:3}
Write the payoff as $(e^{-\int_0^Sr}-c\,e^{-\int_0^Er})^+$, change to the measure with
density $e^{-\int_0^Sr}/P(0,S)$, and price a Black put on $e^{-\int_S^Er}$, lognormal with
mean $P(0,E)/P(0,S)$. Its variance is the bond-price variance at $S$ plus the variance
of the rate's integral over the period, $\approx\sigma^2\delta^3/3$.
\iqlookfor{the measure change and the two variance terms.}
\end{solution}

\begin{solution}{iq:rc:modelling-overnight-rate-products:4}
The overnight rate moves only on meeting days, so an option's uncertainty resolves in
steps: vega and theta concentrate on meetings, and after the last meeting in a period
the period's rate is known. Hedging is by meeting, with futures and their options.
\iqlookfor{step risk and meeting-by-meeting hedging.}
\end{solution}

\begin{solution}{iq:rc:modelling-overnight-rate-products:5}
Discretisation before statistics: whether the simulation grid contains the period's
start and end (a truncated step count shortens the period), the compounding
convention, then convergence in the step size, the curve's forwards used for the
drift, and the discounting; compare with the forward-looking case, which has a closed
form too.
\iqlookfor{bias versus noise, and the grid.}
\end{solution}

\begin{solution}{iq:rc:modelling-overnight-rate-products:6}
It is computed from overnight-rate derivatives; large derivative activity on it would
be circular and thin the underlying market. Official guidance limits its derivative
use to end users hedging term-rate loans; banks hedge those with overnight-rate
instruments and carry the basis.
\iqlookfor{circularity and where the liquidity is.}
\end{solution}
